% !TeX root = ../main.tex
\begin{abstract}
Software architecture descriptions can diverge from the source code they are intended to govern. This paper presents ArchSync, a deterministic architecture-conformance prototype that treats a machine-readable architecture model as an executable contract and reconstructs an observed service graph from TypeScript source code. Stack-specific analyzers detect selected HTTP, PostgreSQL, Redis, and AMQP interactions; a conformance engine then classifies a change as no-impact, rule violation, or architecture evolution and attaches model and file-line-column evidence. A Git-diff gate caches the base graph, incrementally rescans affected components, and maps these classes to PASS, BLOCK, or REVIEW without silently rewriting the approved model. We evaluate the prototype with two controlled datasets and a change-scoped protocol. An Order Platform benchmark contains 20 independently applied patches: nine no-impact changes, seven violations, and four evolutions. Full-graph and changed-graph detection had no false positives or false negatives; ArchSync classified 20/20 cases, matched 7/7 violated-rule sets, and localized 11/11 finding-bearing cases to the expected file and line. Replaying the same patches as Git diffs matched 20/20 merge decisions, changed-file sets, and architecture deltas; every incremental result also matched a separate full scan of the same head repository. The gate hit the baseline cache on 20/20 repeated checks and parsed 57 of 189 TypeScript file instances. On the recorded Windows environment, cold and warm median check latency was 518.51 and 242.62 ms. A separate 40-signal challenge corpus improved from 18 true positives, four false positives, two false negatives, and 16 true negatives in frozen v0.1 to 20, zero, zero, and 20 in provenance-aware v0.2. These are feasibility and regression results, not estimates of general accuracy: the datasets were co-developed with the analyzer, the system is synthetic, timings come from one machine, and only explicit TypeScript/Node.js patterns are covered.
\end{abstract}
